\documentclass[10pt]{exam}
\usepackage[phy]{template-for-exam}

\title{Dilemma}
\author{Rohrbach}
\date{\today}

\begin{document}
\maketitle


\section*{The Problem}

While washing your car, water comes out of the hose at a rate of 1.5 kg/s and a speed of 20 m/s.  How much force does the water impart to your car?

\section*{My Solution (Algebra-based)}

Watch the hose for $\Delta t=1$~s.  The mass that comes out is $m=1.5$~kg.  The water hits the car and comes to a stop, so $\Delta v=-20$~m/s.  Thus,
%
\begin{align}
  F &= \frac{\Delta p}{\Delta t} \,, \\
    &= \frac{m \Delta v}{\Delta t} \, , \\
    &- \SI{-30}{N} \,.
\end{align}
%
Since \SI{-30}{N} is imparted on the water, \SI{30}{N} is imparted on the car.


\section*{My Colleague's Solution (Product Rule)}

By the product rule,
%xw
\begin{align}
  F &= \frac{dp}{dt} \,, \\
    &= \frac{d}{dt}\left(mv\right) \,, \\
    &= m \frac{dv}{dt} + v \frac{dm}{dt} \,. \label{pr}
\end{align}
%
Since $dm/dt=1.5$~kg/s, the second term works out to 30~N.  His argument is that velocity is constant, so $dv/dt=0$.

I'm not convinced that $dv/dt=0$.  But on the other hand, I can't make sense of what $dv/dt$ means in this case.  It's the acceleration of what? Of the infinitesimal bit of water that hits the car in time slice $dt$?  If so, I don't know how I'd calculate it.  Maybe it doesn't matter becase the $m$ is infinitessimal, so that the $m$ is what makes the first term zero in Eq.~(\ref{pr}) instead of the $dv/dt$?

It is all the more confusing because the assumption I made in my solution was that the water comes to a stop when it hits the car.  I could imagine a situation in which this is not true (\emph{i.e.} a Pelton water wheel).  But I don't know how to take that into account using the derivatives.  For example, if we imagine the water from the hose bouncing off of the car elastically, how would that change the solution?  I assume it would double $dv/dt$, but if $dv/dt$ is zero, that doesn't work.  



\end{document}